% vim: tw=78 ai spell spelllang=cs

\documentclass{article}

\usepackage{graphicx}
\usepackage{hyperref}
\usepackage[utf8x]{inputenc}
\usepackage[czech]{babel}
\usepackage[T1]{fontenc}
\usepackage{xspace}

\pagestyle{empty}

\newcommand{\topinfo}[4]{
\begin{flushright}
\today
\end{flushright}
\vspace{1em}
\begin{center}
\Large\textbf{Posudek #1 na #2 práci \\
\normalsize%
\textit{#3:} #4}
\end{center}
}
\newcommand{\botinfo}{
\vspace{3em}
\begin{flushright}
Jiří Slabý
\end{flushright}
}


%obsah strana VIII
%Uvod, 2. veta
%3.odst chaos
%2.1 -- nesrozumitelne
%2.2 odst. 2 konec nesmysl
%obr+zdrojak 2.1 je blbe

\begin{document}
\topinfo{vedoucího}{bakalářskou}{Lukáš Linhart}{Automatické hledání chyb ve
zdrojových kódech pomocí symbolické exekuce}

Dle zadání měla bakalářská práce nastudovat použití symbolické exekuce na
velké projekty a navrhnout a implementovat nástroj, který pomůže při nasazení
exekuce na takové projekty. Text práce je na 18 stranách strukturovaný do 5
kapitol: úvod a motivace, základy nutné pro další text, návrh a implementace,
závěr a manuál. \textbf{Rozsah} práce se pohybuje u samotné spodní hranice
doporučených požadavků na bakalářské práce.

\textbf{Obsahově} by si práce zasloužila větší péči. Je zde mnoho pravopisných
i faktických chyb (např.\ Obrázek a Zdrojový kód 2.1) a trochu chaotického
textu např.\ již v 3.\ odstavci úvodu.  Navíc někde jednotlivé odstavce i
kapitoly nepříliš navazují.  Autor mohl také lépe a rozsáhleji popsat
vlastnosti jednotlivých návrhů včetně nevýhod, čímž by se vzdálil od
minimálního rozsahu.

\textbf{Praktická část} práce je nástroj \textit{dtlu}, v IS MU jako
\texttt{zip} soubor (u textu práce chybí). Dle demonstrace se nástroj tváří
funkčně, dle očekávání. S autorem jsem v průběhu implementace komunikoval a
on pružně reagoval na požadavky, které do práce postupně integroval.

Na závěr konstatuji, že přes uvedené připomínky práce dle mého názoru
\textbf{splňuje zadání a požadavky na bakalářskou práci}.  Doporučuji práci
\textbf{uznat} jako bakalářskou, avšak zvážíme-li rozsah a kvalitu textu
práce, doporučuji hodnocení stupněm C.

\textbf{Nutno podotknout, že tištěná verze se výrazně liší od elektronické v
IS MU. Nechť se student vyjádří také k tomuto a provede nápravu.}

\botinfo
\end{document}
